% =====================================================================
% PROPOSED EDITS -- WHEEL SPEED, INERTIA TARGET, WHEEL GEOMETRY
% From: Pablo.  For: Dejan.  Prepared 28 Jul 2026.
%
% NOTHING APPLIED. These are drop-in replacements for existing content in
% 04_modeling_control.tex, 05_hardware_design.tex and 08_bom.tex.
%
% Basis: electrical/MOTOR_SPEED.md (speed), electrical/WHEEL_130_M6.md
% (wheel), electrical/ENVELOPE_JUSTIFICATION.md (180 mm envelope).
%
% ONE NUMBER CHANGES EVERYTHING: omega_max = 6000 rpm (628 rad/s), replacing
% 8400 rpm in Sections 4-5 and 7000 rpm in the BOM. Both existing figures
% were optimistic; 6000 rpm is a deliberate conservative design value.
% =====================================================================


% ---------------------------------------------------------------------
% EDIT 1 -- 04_modeling_control.tex, tab:jumpup_budget (line ~182)
% Replace the omega_max, h_w and I_w,target rows.
% ---------------------------------------------------------------------
% OLD:
%   $\omega_{\max}$    & Max wheel speed              & $8400$ rpm ($880$ rad/s) \\
%   $h_w$              & Required angular momentum    & $0.277$ kg\,m$^2$/s \\
%   $I_{w,\text{target}}$ & \textbf{Target wheel inertia} & $\mathbf{3.15\times10^{-4}}$ ...
% NEW:

$\omega_{\max}$    & Max wheel speed (design)     & $6000$ rpm ($628$ rad/s) \\
$h_w$              & Required angular momentum    & $0.277$ kg\,m$^2$/s \\
$I_{w,\text{target}}$ & \textbf{Target wheel inertia} & $\mathbf{4.40\times10^{-4}}$ \textbf{kg\,m}$^\mathbf{2}$ \\


% ---------------------------------------------------------------------
% EDIT 2 -- 04_modeling_control.tex, after tab:jumpup_budget
% New paragraph justifying the speed. Currently the document asserts a
% speed with no derivation at all.
% ---------------------------------------------------------------------

\paragraph{Choice of $\omega_{\max}$.} The maximum wheel speed is a design
choice rather than a datasheet figure, and it enters the sizing directly
through $I_{w,\text{target}} = h_w/\omega_{\max}$. At $K_v =
\qty{380}{rpm\per\volt}$ on a nominal \qty{22.2}{\volt} pack the motor's
no-load speed is \qty{8436}{rpm}, but this is unattainable in service: winding
resistance ($2\times\qty{194}{\milli\ohm}$ line-to-line), the reduced phase
voltage available under sinusoidal field-oriented modulation, and the fall of
pack voltage through the discharge curve together place the realistic ceiling
near \qty{6900}{rpm}. The design value is set to \textbf{\qty{6000}{rpm}
(\qty{628}{\radian\per\second})}, below even that figure. At this speed the
back-EMF is \qty{15.8}{\volt}, only \qty{71}{\percent} of the nominal bus, so
the wheel reaches its commanded speed across the whole usable discharge range
rather than only on a freshly charged pack, and margin remains for motor
tolerance, bearing and aerodynamic drag, and the residual uncertainty in the
mass budget. The cost is a \qty{40}{\percent} larger inertia requirement than
a top-of-range assumption would give, which the wheel of
Section~\ref{sec:rw_sizing} meets with \qty{2.4}{\percent} margin. The
corresponding electrical frequency, \qty{1200}{\hertz} at 12 pole pairs, sits
well inside the driver's capability.


% ---------------------------------------------------------------------
% EDIT 3 -- 05_hardware_design.tex, tab:inertia_cases (line ~78)
% Caption and both data rows. Recomputed at 6000 rpm.
% ---------------------------------------------------------------------
% OLD caption: {Required per-wheel inertia by contact case, at $\omega_{\max}=8400$~rpm.}
% NEW:

\caption{Required per-wheel inertia by contact case, at $\omega_{\max}=6000$~rpm.}
% ... and the two data rows become:
Edge                       & 0.251 & 3.99 \\
\textbf{Corner (critical)} & \bfseries 0.277 & \bfseries 4.40 \\

% NOTE: the surrounding prose says the corner case requires
% "$3.15\times10^{-4}$ kg m^2 against $2.85\times10^{-4}$ kg m^2 for the edge
% ---about $10\%$ more". The 10% relation is unchanged (10.2%), but both
% absolute numbers move: 4.40 and 3.99 respectively.


% ---------------------------------------------------------------------
% EDIT 4 -- 05_hardware_design.tex, sec:rw_sizing
% The selected wheel changes from 120 mm/21 nuts to 130 mm/18 nuts, and the
% nut changes from a mis-specified part to a real M6. Replaces the
% paragraph beginning "The structure is printed in PET-CF".
% ---------------------------------------------------------------------

The structure is printed in PET-CF ($\rho = 1290$ kg/m$^3$): a
$130$~mm outer-diameter ring of $20$~mm radial width and $12$~mm thickness,
carried by three radial spokes $10$~mm wide of the same thickness. The hub and
motor interface are neglected in the inertia budget. The ballast is standard
M6 steel hex nuts (ISO 4032: $10$~mm across flats, $5.2$~mm thick,
$\rho = 7850$ kg/m$^3$, measured mass $2.5$~g), seated in blind hexagonal
pockets machined to $5.2$~mm depth at the mid-width of the ring, leaving a
$6.8$~mm plastic floor beneath each nut. Accounting for the plastic each
pocket displaces, the \emph{net} contribution per installed nut is $+1.87$~g
of mass and $+5.70\times10^{-6}$ kg\,m$^2$ of inertia at the ring radius.

The selected design is $\mathbf{OD = 130}$~\textbf{mm with 18 nuts}: it reaches
$4.51\times10^{-4}$ kg\,m$^2$ against the $4.40\times10^{-4}$ kg\,m$^2$ target
($102.4\%$), leaves $7.2$~mm of plastic between adjacent pockets and at least
$4.0$~mm of radial wall in either nut orientation, and keeps the three-wheel
set to $\sim\!485$~g. The printed ring supplies $74\%$ of the total inertia
and the ballast $23\%$, so the nuts act as a trim about a structure that
nearly meets the target unaided, and the wheel can be adjusted in either
direction after fabrication.

\paragraph{Pocket placement.} The pockets cannot be spaced uniformly. A spoke
occupies $\pm6.4^\circ$ where it meets the ring and a pocket $\pm6.3^\circ$ at
the ballast radius, so a pocket centre must clear a spoke centre by
$12.7^\circ$; with three spokes, no uniform pitch at any usable nut count
satisfies this. The eighteen pockets are therefore distributed as six per
$120^\circ$ sector between spokes, spread over the $94.7^\circ$ of usable arc
at a $18.9^\circ$ pitch. Three-fold symmetry, and hence balance, is preserved,
and the arrangement admits at most seven per sector should additional trim be
required.

\paragraph{Nut retention.} At \qty{6000}{rpm} each nut experiences
\qty{54}{\newton} of centrifugal load at the \qty{55}{\milli\metre} ballast
radius. The blind-pocket floor carries this in bending, but positive retention
is required: a cover plate over the pocket mouths is preferred to adhesive,
since it preserves the post-fabrication adjustability that motivates the
ballast scheme.


% ---------------------------------------------------------------------
% EDIT 5 -- 05_hardware_design.tex, sec:massbudget volume constraint
% Replaces the "$180$~mm cube envelope" sentence with a justified version.
% ---------------------------------------------------------------------

\paragraph{Volume constraint.} Every subsystem---the three orthogonal
reaction-wheel/motor modules, the battery, and the electronics stack---must fit
inside the $180$~mm cube envelope. This envelope is not arbitrary. Sizing the
wheels to the conservative \qty{6000}{rpm} of Section~\ref{sec:jumpup_target}
requires $4.40\times10^{-4}$ kg\,m$^2$ per wheel; within a $150$~mm envelope
the largest wheel that clears the frame, motor and contact features is about
$120$~mm in diameter, yielding $2.63\times10^{-4}$ kg\,m$^2$---$17\%$ short,
and not recoverable by ballast, which necessarily sits inboard of the ring it
displaces. The $180$~mm envelope admits the $130$~mm wheel of
Section~\ref{sec:rw_sizing}. Packaging independently demands it: the
$150\times90$~mm power and signal board does not fit within a $150$~mm cube at
any wall thickness. The binding dimension remains the reaction wheel, whose
outer diameter caps the radius at which ballast mass can be placed.


% ---------------------------------------------------------------------
% EDIT 6 -- 08_bom.tex, tab:bom-derived
% Caption is wrong on both cube figures, and three rows are wrong.
% ---------------------------------------------------------------------
% OLD caption: "...assuming a \qty{150}{\milli\metre}, \qty{2}{\kilo\gram} cube"
% NEW caption:

\caption{Derived actuation parameters (computed from the supplied BOM; order-of-magnitude, assuming a \qty{180}{\milli\metre}, \qty{1.45}{\kilo\gram} cube).}

% Replacement rows:

Usable wheel speed (design) & conservative value, see Sec.~\ref{sec:jumpup_target}; \qty{8436}{\rpm} no-load nominal & \qty{6000}{\rpm} (\qty{628}{\radian\per\second}) & Design choice \\
Electrical frequency & 12 pole pairs at \qty{6000}{\rpm} & \qty{1200}{\hertz}, inside driver capability & Derived \\
Required jump-up momentum & corner case, $\eta=1.05$, $\beta=1.34$ & \qty{0.277}{\kilogram\metre\squared\per\second} & Derived \\
Target wheel inertia & $h_w/\omega_{\max}$ at \qty{628}{\radian\per\second} (corner governs) & \qty{4.40e-4}{\kilogram\metre\squared} & Derived \\
Stored kinetic energy / wheel & $\tfrac12 J\omega^2$ at \qty{628}{\radian\per\second} & \qty{89}{\joule} (\qty{267}{\joule} for 3) & Derived \\


% ---------------------------------------------------------------------
% EDIT 7 -- 08_bom.tex, tab:bom-gaps
% The rim row procures the wrong wheel entirely (AUDIT.md M-8).
% ---------------------------------------------------------------------
% OLD:
%   Reaction-wheel rims & 3 & ... laser-/water-cut steel ring, ~120 mm OD, 3-4 mm thick & H \\
% NEW (the wheel is printed, not cut; what must be bought is the ballast):

M6 steel hex nuts (ISO 4032) & 60+ & Wheel ballast: 18 per wheel $\times$ 3, plus spares for post-fabrication trim & L \\

% and delete the steel-rim row, since the wheel is printed in PET-CF and is
% covered by the existing filament line.
% NOTE: this also removes one of the two items the BOM introduction calls the
% leading long-lead schedule risk. That paragraph needs revisiting.


% ---------------------------------------------------------------------
% NOT PATCHED -- still open, listed in electrical/AUDIT.md
%
% 1. Total cube mass M = 1.45 kg is contradicted by its own budget: wheels
%    (485 g at the new geometry) + battery (600 g) + 3 motors (204 g) is
%    1289 g, i.e. 89% of target, with no frame, electronics or fasteners.
%    h_w scales with M, so if M grows the inertia target grows with it and
%    this wheel stops closing. THIS IS NOW THE LARGEST OPEN NUMBER.
%
% 2. tau_b = 5.0 N m per wheel has no stated source and the supplied MG92B
%    servo produces 0.29 N m. beta depends on tau_b directly.
%
% 3. 08_bom.tex still contains unresolved merge-conflict markers and does
%    not compile.
%
% 4. Motors/drivers/encoders: 3 confirmed, no spare. The Table 1 footnote
%    claiming four are procured, and the Table 3 spare row, are both wrong.
% ---------------------------------------------------------------------
